\section{Scope, conventions, and the status of the results}
\label{scope:sec}

This continuation develops two exact extensions of the ProveIt programme.
First, it answers the archived question about an all-index half-twist
Stieltjes multiplication law with its endpoint terms. The twist changes
the convolution identity from \(\delta_0-1\) to \(\delta_0\).
Second, it determines the principal Laurent coefficients and finite
parts of a shifted harmonic Dirichlet family at every nonpositive
integer, with reflection, rational multiplication, and antiderivative
formulas for those coefficients.

Both developments require a precise normalization at a moving pole.
The periodic correction is supported at the endpoint of a circle.
The harmonic correction is a Laurent residue, which survives when
finite parts are differentiated. Their common role explains why a
formula at regular parameters cannot simply be specialized or
differentiated at a singular parameter.

\subsection{The inspected baseline}

We froze the repository at \pinned, whose commit is dated
10 October 2026, 21:44:12 UTC. The audit used the complete set of
72 chapter sources and the eight non-PDF manuscript support files.
The pending archive
\path{ProveIt_Polylogarithm_Stieltjes_Continuation_2026-10-10.zip}
was matched to its Git blob and its mathematical source was inspected.
Five relevant archived report sources were also checked.
The machine-readable source manifest records exact paths and hashes.
These statements describe the inspected set; they do not assert that
every report or historical archive in the repository was reviewed.
Manuscript and archival status are kept separate
\cite{RepoManuscript,RepoIncoming}.

The archived \emph{Convolution Algebra} report explicitly asks for a
nonintegral-twist normal form and, as a first target, the all-index
half-twist multiplication table with point-supported counterterms
\cite[final research section]{RepoConvolutionAlgebra}.
Section~\ref{tw:sec:main} supplies this normalization, multiplication
law, and covariant derivative/primitive calculus. It gives a
rational-twist conversion to Hurwitz and Stieltjes coordinates, a
convergent half-twist integral, and the exact zero-frequency correction
needed in the limit as the twist tends to zero.

The same archive and two further continuations already contain the
untwisted periodic Bell algebra, polygamma contact terms, normalized
primitive tower, and reflected polylogarithm kernels
\cite{RepoConvolutionAlgebra,RepoConvolutionCalculus,RepoShiftedJets}.
Those results are consolidated with independent complete proofs here;
they are not new theorems relative to the pinned repository.
The incoming Stieltjes continuation supplies additional ordinary
translation and integration formulas \cite{RepoIncoming}.
Earlier alternating-harmonic reports likewise already use beta
representations, cyclotomic projections, and scalar Mellin subtraction
\cite{RepoAlternating,RepoHarmonicGeometry}.

The source's \(S_2\) and \(S_4\) identities already have exact proofs.
The frozen \(S_6\) target and the newer \(S_8\) candidate remain
conjectural. Our proofs do not evaluate either of them.
The half-twist target answered here is an explicit research question;
it is distinct from those two special-value conjectures.

\subsection{Classical inputs and the contribution of this article}

The Fourier coefficients of Hurwitz zeta, Gamma--beta identities, and
\[
 \zeta'(0,x)=\log\Gamma(x)-\frac12\log(2\pi)
\]
are classical \cite{DLMFHurwitz,DLMFBeta}.
Normalized Hurwitz convolution is established; Sun states an
ordinary-integral form in his study of invariant functions
\cite[Remark 3.2 in the arXiv version]{per:Sun}.
Periodic complex-power distributions provide the natural background
\cite{per:GaySebbar}. Kim's Hurwitz-transform continuation is a recent
antecedent for scalar test-function subtraction \cite{Kim2026}.

Young's parametric Euler sums supply the beta-derivative evaluation
framework \cite{Young2026}. His height-one work and unshifted Laurent
study are directly relevant \cite{Young2024,Young2023}; the harmonic
Stieltjes literature concerns related, but not identically shifted,
families \cite{HarmonicStieltjes2023}. The full text of Young's 2023
paper was unavailable. Its inspected publisher abstract already rules
out a broad priority claim for unshifted Laurent expansions.

The additions for this continuation are the fully normalized twisted
extension and the joint coefficient theorem for the outer-shifted
harmonic family, with its finite polynomial calculus.
Beta differentiation, Mellin continuation, and the unshifted height-one
theory are treated as antecedents. Some consequences may have
equivalent formulations elsewhere. A result described as \emph{proved
here} has a complete analytic proof in this article; a result described
as an \emph{addition to the inspected source} was not located in that
source set. Neither phrase asserts universal literature priority.
Proof and provenance statuses are recorded separately.

\begin{center}
\begin{tabularx}{\textwidth}{@{}p{0.25\textwidth}X@{}}
\toprule
Result or ingredient & Status and use\\
\midrule
Untwisted periodic calculus &
Existing repository foundation, independently checked; all endpoint
terms and integration constants retained.\\[3pt]
Twisted Stieltjes calculus &
Answers the archived half-twist target, for every real nonintegral
twist, with the singular normalization included.\\[3pt]
Ordinary half-twist integral &
A convergent consequence with a midpoint evaluation in
\(\log\Gamma(1/4)\), Euler's constant, and logarithms.\\[3pt]
Harmonic beta evaluations &
Established antecedent connecting harmonic polynomials, Gamma
derivatives, and positive integer values.\\[3pt]
Shifted harmonic Laurent data &
Complete principal and finite parts; residue-aware differentiation,
reflection, rational multiplication, and repeated primitives.\\[3pt]
Numerical checks &
Independent diagnostics; none replaces an analytic proof.\\
\bottomrule
\end{tabularx}
\end{center}

\subsection{Notation and normalizations}

Order derivatives always concern the first argument:
\[
 \zeta^{[j]}(s,a)=\partial_s^j\zeta(s,a),\qquad
 \rLi_s^{[j]}(z)=\partial_s^j\rLi_s(z).
\]
By contrast, \(\psi^{(k)}(a)\) differentiates the argument \(a\), and
\begin{equation}\label{scope:stieltjes}
 \zeta(1+w,a)=\frac1w+
       \sum_{n\geq0}\frac{(-1)^n}{n!}\gamma_n(a)w^n.
\end{equation}
Thus \(\gamma_0(a)=-\psi(a)\), and \(\gamma=\gamma_0(1)\).
Finite parts are Laurent constant terms, unless a displayed cutoff
definition specifies the equivalent integral normalization.

On \(\bbT=\bbR/\bbZ\), write
\[
 \widehat T(n)=\langle T,\rme^{-2\pi\numi n x}\rangle,\qquad
 \Delta=\delta_0-1.
\]
The distribution \(\Delta\) is the identity for convolution on the
subspace of distributions with zero mean. The full algebra has identity
\(\delta_0\). Reflection means
\(\check T(x)=T(-x)\), and \(D\) denotes the distributional argument
derivative. A convolution is always the compact-circle convolution.
It is defined for any pair of distributions on the circle; it must
not be confused with a pointwise product of singular distributions.
All logarithms of \(2\pi\numi n\), for \(n\ne0\), mean
\[
 \log(2\pi\numi n)=\log(2\pi|n|)
                 +\frac{\pi\numi}{2}\operatorname{sgn}(n).
\]
This branch convention is part of the identities.

For finite harmonic sums we use
\[
 H_N^{(j)}=\sum_{k=1}^N k^{-j},\qquad
 H_N(\{1\}^r)=
 \sum_{N\geq n_1>\cdots>n_r\geq1}\frac1{n_1\cdots n_r},
 \qquad H_N(\{1\}^0)=1.
\]
Their generating polynomial is
\[
 \sum_{r\geq0}H_N(\{1\}^r)u^r
      =\prod_{k=1}^N(1+u/k).
\]
Thus all depths concern explicit polynomials in generalized harmonic
numbers:
\[
 H_N(\{1\}^2)=\frac{H_N^2-H_N^{(2)}}2,\qquad
 H_N(\{1\}^3)=\frac{H_N^3-3H_NH_N^{(2)}+2H_N^{(3)}}6.
\]
